\section{Key-Release Delivery for Digital Goods (Specified)}
\label{sec:delivery}

Theorem~\ref{prop:loss} stops at capture: a vendor that captures and delivers nothing is paid, within the limits. This section specifies one
protocol that extends the protection to delivery, for one kind of good, with both parties' protection and the assumptions it needs. The
dispute predicate is built and measured; the vault's order states, the bond and the window are not built, and Section~\ref{sec:eval} does not
evaluate them.

\paragraph{Scope.} A \emph{digital good} here is a byte string $M$ of at most 4 KiB whose hash $h_m = H(M)$ the vendor commits to before payment.
That covers a quote, a model's answer or a record computed before the 402 is sent; larger goods need a chunked commitment, which we leave to future work.
The protocol decides whether the vendor released a key that opens a committed ciphertext to the committed bytes. It does not decide whether those bytes
are what the buyer wanted; a named data source that signs $h_m$ is the protocol's only way to give them meaning, and quality no such signature captures is
out of scope.

\subsection{The protocol}

\paragraph{Offer.} The vendor draws a fresh key $k$, encrypts $M$ with a keccak keystream (block $i$ of the ciphertext is block $i$ of $M$ XOR
$\mathrm{keccak}(k, i)$, 32-byte blocks), and answers the buyer's request with $c$ and an EIP-712 offer $\sigma_V$ over (resource, asset, payee, price,
$h = H(k)$, $h_c = H(c)$, $h_m$, a fresh order identifier, the buyer named in the request, the vault, the last acceptance time $T_a$ and the delivery
deadline $T$), signed by the key the vendor registered with its bond. Fresh $k$ per order is required: the keystream is not secure under key reuse.

\paragraph{Acceptance.} The buyer's client checks $H(c) = h_c$ and keeps $c$; a buyer that does not hold $c$ must not accept. The relay submits the request
with $\sigma_V$. The vault verifies the signature against the vendor's registered key, that the order identifier is unused and names this buyer, that the
price is the signed price, and that $t < T_a$, with $T - T_a$ at least the agreed inclusion delay $\delta$. The request then goes through the same encrypted
check and debit as any other. On approval the payment is funded into the escrow and locked against $(h, h_c, h_m, T)$, and the vendor's penalty $\pi$ is
reserved from the free balance of its bond in the same transaction, which reverts if too little is free. An order exists only once funded and locked.
Acceptance means this atomic funded transition, not submission of the request: funding rechecks the signature, unused order identifier and
$t < T_a$. Approval that arrives too late creates no order and no penalty liability; its unused reservation can be reclaimed under the base policy.
The order identifier becomes used only at funding. Thus confidential evaluation cannot shorten the vendor's promised response margin.

\paragraph{Release and capture.} The vendor calls \code{claim}$(k)$ before $T$. The vault checks $H(k) = h$, records $k$, and starts a dispute window $W$;
publishing $k$ enables decryption. After $W$ without a valid dispute the escrow captures to the vendor and the penalty reservation is released. No valid
$k$ by $T$: anyone may void the hold, the amount returns to the budget as in a reclaim, and the penalty moves from the bond to the buyer. The trigger is
the absence of $k$ on chain, never the buyer's claim. This needs an extended escrow, not the deployed auth-capture escrow unchanged: a funded order has
exactly one terminal outcome (capture, no-key refund or successful-dispute refund), and neither payer reclaim nor goal closure can bypass its window.
The capture deadline must extend beyond $T + W + \delta$, and a timely release must remain claimable through that window. A refund atomically returns
tokens to the payer, credits the encrypted budget and restores one call exactly once; capture does neither. Neither outcome restores an attempt.
An order still awaiting its window stays reserved in $R$ and $L$ of Section~\ref{sec:model}. These are requirements for the unbuilt integration.

\paragraph{Dispute.} Within $W$, anyone holding $c$ may submit it with a dispute fee $f$. The vault checks $H(c) = h_c$ and $H(k) = h$, decrypts, and
compares $H(\mathrm{Dec}_k(c))$ with $h_m$. A mismatch voids the hold, returns the amount to the payer and the budget, pays the penalty and returns $f$; a
match pays $f$ to the vendor. A ciphertext whose hash does not match $h_c$ is invalid evidence and cannot refund the order. Our implementation of this predicate (\code{KeyReleaseVerifier}) uses 86{,}894 gas of execution for a 4{,}096-byte
ciphertext, plus about 65{,}000 for the calldata. The predicate and calldata therefore fit under 200{,}000 gas; this excludes the unbuilt order,
bond and refund logic. The fee $f$ must cover the complete dispute's gas cost at an agreed gas price, which has not yet been measured.

\subsection{What each party is protected from}

\paragraph{Buyer.} (D1) The vendor is captured to only if it published, before $T$, a $k$ with $H(k) = h$, and no valid dispute followed within $W$.
(D2) Whenever $k$ is published, a buyer that accepted holds $c$ with $H(c) = h_c$, because it checked before accepting, and can decrypt.
(D3) A vendor that does not publish $k$ by $T$, on an order accepted by $T_a$, loses $\pi$ to the buyer, decided by the chain alone.
(D4) If $\mathrm{Dec}_k(c)$ is not the committed $M$, a dispute within $W$ refunds the price and pays $\pi$. D1 follows from the hash lock and the
window, D2 from the acceptance check, D3 from the atomic reservation and the on-chain trigger, D4 from the commitments and the price held through $W$.

\paragraph{Vendor.} (V1) Once an order is locked the buyer cannot withdraw before $T$: the funds are in the escrow under the vault's control, and nothing the
buyer or its agent can do changes the order's state. (V2) A vendor whose \code{claim} is included before $T$ is captured to unless a dispute proves the
predicate false. A false dispute fails the vault's own recomputation and pays the vendor $f$. (V3) A buyer cannot lock an order so late that the vendor
has no time to claim, because acceptance ends at $T_a \le T - \delta$.

\paragraph{Key disclosure without settlement.} A \code{claim} discloses $k$ when it is broadcast, before it settles. It could then revert, be
delayed past $T$, or be included in a block that is later dropped, leaving a buyer with the good and the vendor unpaid. A claim can revert only if it is late or the key is
wrong, since after acceptance no other party can change the order (V1). A dropped block restores the order to locked with $k$ already public; the vendor
resubmits the same claim, which is valid until $T$ on the canonical chain. The vendor is therefore exposed only if its claim, submitted by $T - \delta$, is
not included by $T$ on the canonical chain, which is the same inclusion assumption D3 makes in the buyer's favour. A vendor wanting to avoid public
disclosure before inclusion can submit through a private transaction channel; we do not rely on one.

\subsection{Assumptions and what remains open}

\emph{Inclusion:} a transaction submitted $\delta$ before a deadline is included before it, on the canonical chain. \emph{An honest disputer:} under A2
the buyer's agent may decline to dispute. D4 then needs another party holding $c$, such as the user's console or a watcher the user runs; without one,
the loss is still bounded by Theorem~\ref{prop:loss}, and D1 to D3 do not depend on disputes. \emph{Privacy:} $k$ is public, so $M$ stays private only as
long as $c$ does. A dispute publishes $c$ and therefore $M$; a buyer who disputes gives up the good's confidentiality. \emph{Collusion:} the allowlist,
not the offer signature, limits which vendors can be paid; a signed offer stops an orchestrator altering an honest quote and does nothing against an
allowlisted vendor that signs an inflated one. \emph{Adoption:} a vendor must accept the vault as the escrow's operator and register a bond.

The intended bilateral protection is conditional on the extended escrow and accounting rules above, the inclusion assumption and, for D4, someone
disputing. Only the predicate has been implemented; no end-to-end fair-exchange guarantee has been established. The protocol cannot judge goods
whose value is not fixed by a commitment made before payment.
Held services, such as a reserved GPU slot, need a signed hold and a different proof of delivery, and remain future work (Section~\ref{sec:limits}).

\begin{figure*}[t]
\centering
\resizebox{0.88\linewidth}{!}{\begin{tikzpicture}[x=1cm,y=-0.62cm,
    msg/.style={-{Stealth[length=2mm]},thick},
    rmsg/.style={-{Stealth[length=2mm]},thick,dashed},
    mlab/.style={font=\scriptsize,fill=white,inner sep=1pt,align=center},
    act/.style={draw,thick,fill=white,font=\scriptsize\bfseries,align=center,minimum width=2.2cm,minimum height=0.75cm,rounded corners=2pt}]
  \def\xB{0}\def\xE{6.2}\def\xV{12.4}
  \node[act,draw=untrusted] at (\xB,0) {Buyer\\(console, agent)};
  \node[act,draw=chain]     at (\xE,0) {Vault\\and escrow};
  \node[act,draw=untrusted] at (\xV,0) {Vendor\\(bonded)};
  \foreach \x in {\xB,\xE,\xV} \draw[gray!60,thick] (\x,0.8) -- (\x,14.0);
  \draw[msg]  (\xB,1.6) -- node[mlab,above]{GET /resource} (\xV,1.6);
  \draw[rmsg] (\xV,2.8) -- node[mlab,above]{402: $c$, offer $\sigma_V$ over price, $h{=}H(k)$, $h_c{=}H(c)$, $h_m{=}H(M)$, order id, buyer, $T_a$, $T$} (\xB,2.8);
  \node[mlab,anchor=west] at (\xB+0.1,3.8) {check $H(c) = h_c$; keep $c$};
  \draw[msg]  (\xB,4.9) -- node[mlab,above]{request with $\sigma_V$, before $T_a$} (\xE,4.9);
  \node[mlab,anchor=west,text width=5.4cm] at (\xE+0.15,6.0) {verify offer; encrypted check and debit; fund and lock against $(h, h_c, h_m, T)$; reserve penalty $\pi$ from the bond};
  \draw[msg]  (\xV,7.4) -- node[mlab,above]{claim($k$) before $T$} (\xE,7.4);
  \node[mlab,anchor=west,text width=5.4cm] at (\xE+0.15,8.5) {check $H(k) = h$; record $k$; open window $W$};
  \draw[rmsg] (\xE,9.6) -- node[mlab,above]{$k$ public: delivery} (\xB,9.6);
  \node[mlab,anchor=west] at (\xB+0.1,10.5) {$M = c \oplus \mathrm{keccak}(k, i)$};
  \node[mlab,anchor=west,text width=5.4cm] at (\xE+0.15,11.3) {after $W$: capture to the vendor, release $\pi$};
  \node[draw,dashed,rounded corners,font=\scriptsize,align=center,fill=orange!8,text width=12.6cm,anchor=north] at (\xE,12.3)
    {No valid $k$ by $T$: hold voided, amount back to the budget, $\pi$ to the buyer. Within $W$, anyone holding $c$ may dispute with fee $f$: the vault checks $H(c){=}h_c$, decrypts and compares with $h_m$. Mismatch: price refunded, $\pi$ and $f$ to the buyer. Match: $f$ to the vendor.};
\end{tikzpicture}
}
\caption{The key-release protocol. Only the dispute predicate is built. Publishing $k$ and starting the window are one transaction; every other
outcome is decided by what is or is not on chain.}
\label{fig:exchange}
\end{figure*}
